\chapter{Морфизм \texorpdfstring{$K_{43}$}{K43}-отклика в плотность выхода}
\label{ch:v10-breathing-k43-output-density-morphism}

\section{Точный знак ориентированного отклика}

Положим $x=e^\zeta>0$. Следовой отклик ориентированной входящей
самоэнергии переписывается без логарифмов:
\begin{equation}
 \mathcal A(q,x)
 =\frac{x}{1+q^2+x}-\frac{x}{1+x}
 =-\frac{xq^2}{(1+x)(1+q^2+x)}.
 \label{eq:v10-k43-output-trace-response}
\end{equation}
Для $q>0$ он строго отрицателен. Из двух линейных знаковых отображений
$\mathcal A\mapsto\pm\mathcal A$ условию положительного выхода удовлетворяет
только
\begin{equation}
 R_{\rm out}(q,x):=-\mathcal A(q,x)
 =\frac{xq^2}{(1+x)(1+q^2+x)}>0.
 \label{eq:v10-k43-output-fraction}
\end{equation}
Кроме того,
\begin{equation}
 1-R_{\rm out}
 =\frac{q^2+(x+1)^2}{(x+1)(q^2+x+1)}>0,
 \label{eq:v10-k43-output-subunit}
\end{equation}
поэтому $R_{\rm out}$ является допустимой субединичной потерей на одну
ячейку. В канонической точке $q=x=1$ получаем
\begin{equation}
 R_{\rm out}(1,1)=\frac16.
 \label{eq:v10-k43-output-witness}
\end{equation}

\section{Перевод из потери на ячейку в плотность}

Используем уже выбранную геометрическую меру
\begin{equation}
 v_{\rm cell}=\frac{\beta_E^2}{4\alpha^2m^2}.
\end{equation}
Единственный локальный перевод без нового размерного коэффициента равен
\begin{equation}
 d_{\rm out}
 :=\frac{R_{\rm out}}{v_{\rm cell}}
 =\frac{4\alpha^2R_{\rm out}}{\beta_E^2}m^2
 =\varepsilon_{43}m^2,
 \qquad
 \varepsilon_{43}
 =\frac{4\alpha^2R_{\rm out}}{\beta_E^2}.
 \label{eq:v10-k43-output-density}
\end{equation}
При $q=x=1$ условный коэффициент равен
\begin{equation}
 \varepsilon_{43}^{*}
 =\frac{2\alpha^2}{3\beta_E^2}.
 \label{eq:v10-k43-output-epsilon-witness}
\end{equation}
Таким образом, число $1/6$ не является самим коэффициентом плотности:
оно является выходной долей на ячейку, которую ещё необходимо умножить на
$1/v_{\rm cell}$.

\section{Баланс со сквозным притоком}

Для
\begin{equation}
 d_{\rm in}=\frac{n_{\rm flow}\log2}{v_{\rm cell}}
\end{equation}
разность принимает вид
\begin{equation}
 d_{\rm in}-d_{\rm out}
 =\frac{4\alpha^2m^2}{\beta_E^2}
 \left(n_{\rm flow}\log2-R_{\rm out}\right).
 \label{eq:v10-k43-output-balance}
\end{equation}
Ненулевая стационарная ветвь выбирает слепое условие
\begin{equation}
 \boxed{n_{\rm flow}\log2=R_{\rm out}},
 \label{eq:v10-k43-output-blind-balance}
\end{equation}
но общий множитель $m^2$ вновь сокращается.

\section{Условный родитель и физическая граница}

Нормированные отношения знака $u_R$, плотности $u_d$ и баланса $u_B$
собираются в
\begin{equation}
 \mathcal P_{43\to{\rm out}}
 =\frac12(u_R-1)^2
 +\frac12(u_d-u_R)^2
 +\frac12(u_B-u_d)^2.
 \label{eq:v10-k43-output-parent}
\end{equation}
Его гессиан равен
\begin{equation}
 \begin{pmatrix}2&-1&0\\-1&2&-1\\0&-1&1\end{pmatrix},
\end{equation}
имеет ранг $3$, определитель $1$ и ведущие миноры $(2,3,1)$.
Размерная карта сохраняет ранг/ядро $3/1$ и образующий
$(-1,-1,2,-2)$.

\begin{theorem}[Единственный положительный знаковый морфизм]
Для ориентированного $K_{43}$-отклика при $q,x>0$ прямое отображение
$R=\mathcal A$ отрицательно, тогда как $R=-\mathcal A$ строго лежит между
$0$ и $1$. Деление на $v_{\rm cell}$ даёт положительную плотность
$d_{\rm out}=\varepsilon_{43}m^2$.
\end{theorem}

\begin{nogocriterion}[Положительный отклик ещё не является физическим каналом]
Положительность и субединичность допускают $R_{\rm out}$ как выходную долю,
но не доказывают существование KMS-совместимого полностью положительного
канала или GKSL-генератора, поток которого равен $R_{\rm out}$. Кроме того,
$\alpha$, $\beta_E$ и абсолютный масштаб остаются невыведенными.
\end{nogocriterion}

\begin{protocol}[Статус морфизма выхода]\begin{enumerate}
\item унаследованные ингредиенты: \textbf{$4/4$};
\item условная архитектура: \textbf{$10/10$};
\item условное алгебраическое замыкание: \textbf{$8/8$};
\item положительный субединичный морфизм: \textbf{условно $1/1$};
\item физический открытый канал: \textbf{$0/1$};
\item происхождение коэффициентов и масштаба: \textbf{$0/2$};
\item абсолютный масштаб: \textbf{$0/1$};
\item следующий гейт: \textbf{KMS-совместимый выходной канал}.
\end{enumerate}\end{protocol}

Машинный сертификат сохранён в
\path{s2t_v10_cell_birth_four_volume_induced_newton_breathing_anomaly_k43_output_density_morphism_origin_gate_results.json}.